\chapter{KÝ HIỆU VÀ SỐ CHIỀU CỦA MÔ HÌNH}
\label{app:ky-hieu}

\section{Danh mục ký hiệu}
\label{sec:danh-muc-ky-hieu}

Bảng~\ref{tab:a-ky-hieu} tổng hợp các ký hiệu được sử dụng thường xuyên trong luận văn.
Các ký hiệu được giữ nhất quán giữa mô hình toán học và mã nguồn để thuận tiện cho việc
đối chiếu.

\begin{longtable}{@{}l p{11.2cm}@{}}
\caption{Danh mục các ký hiệu chính}
\label{tab:a-ky-hieu}\\
\toprule
\textbf{Ký hiệu} & \textbf{Ý nghĩa vật lý hoặc thuật toán} \\
\midrule
\endfirsthead
\multicolumn{2}{@{}l}{\textit{(tiếp theo trang trước)}}\\
\toprule
\textbf{Ký hiệu} & \textbf{Ý nghĩa vật lý hoặc thuật toán} \\
\midrule
\endhead
\bottomrule
\endfoot
$K$ & Số thiết bị người dùng trong hệ thống. \\
$N$ & Số phần tử của STAR--RIS. \\
$h_{d,k}$ & Hệ số kênh trực tiếp từ trạm gốc đến UE $k$. \\
$\mathbf{g}$ & Véc-tơ kênh từ trạm gốc đến STAR--RIS. \\
$\mathbf{h}_{r,k}$ & Véc-tơ kênh từ STAR--RIS đến UE $k$. \\
$\phi^T_n,\ \phi^R_n$ & Hệ số phức của phần tử STAR--RIS thứ $n$ ở phía truyền và phía phản xạ. \\
$\beta^T_n,\ \beta^R_n$ & Tỷ lệ năng lượng của phần tử thứ $n$ dành cho phía truyền và phía phản xạ. \\
$\theta^T_n,\ \theta^R_n$ & Pha điều khiển của phần tử thứ $n$ ở phía truyền và phía phản xạ. \\
$h^{\mathrm{eff}}_k$ & Kênh hiệu dụng từ trạm gốc đến UE $k$ sau khi cộng đường trực tiếp và đường qua STAR--RIS. \\
$g_k = |h^{\mathrm{eff}}_k|^2$ & Độ lợi công suất của kênh hiệu dụng UE $k$. \\
$p_c$ & Công suất phát dành cho luồng chung của RSMA. \\
$p_k$ & Công suất phát dành cho luồng riêng của UE $k$. \\
$P_{\max}$ & Tổng công suất phát khả dụng của trạm gốc. \\
$P_{\mathrm{priv}}$ & Tổng công suất của tất cả luồng riêng. \\
$\eta_k$ & Tỷ lệ tốc độ luồng chung phân bổ cho UE $k$. \\
$C_k$ & Phần tốc độ luồng chung thực tế cấp cho UE $k$. \\
$\gamma_{c,k}$ & SINR khi UE $k$ giải mã luồng chung. \\
$\gamma_{p,k}$ & SINR khi UE $k$ giải mã luồng riêng của mình sau SIC. \\
$R_c$ & Tốc độ luồng chung mà tất cả UE đều có thể giải mã. \\
$R_{p,k}$ & Tốc độ luồng riêng của UE $k$. \\
$R_k$ & Tổng tốc độ của UE $k$, gồm phần chung được cấp và phần riêng. \\
$R_{\mathrm{sum}}$ & Tổng tốc độ của tất cả UE. \\
$R_{\min}$ & Ngưỡng tốc độ tối thiểu dùng để đánh giá QoS. \\
$V$ & Tổng phần thiếu hụt QoS của các UE. \\
$V_2$ & Tổng bình phương phần thiếu hụt QoS của các UE. \\
$\sigma^2$ & Công suất tạp âm. \\
$s$ & Véc-tơ trạng thái đưa vào mạng học tăng cường sâu. \\
$a$ & Véc-tơ hành động chuẩn hóa do mạng sinh ra. \\
$d_s$ & Số chiều trạng thái. \\
$d_a$ & Số chiều hành động. \\
$\rho_c$ & Tỷ lệ tổng công suất dành cho luồng chung sau bước giải mã hành động. \\
$w_k$ & Trọng số chia phần công suất riêng còn lại cho UE $k$. \\
$\kappa$ & Hệ số độ sắc của hàm softmax dùng cho phân bổ công suất riêng. \\
$\bar{\theta}_n$ & Pha tham chiếu giải tích của phần tử STAR--RIS thứ $n$. \\
$\lambda_t$ & Hệ số phạt QoS thích nghi ở bước huấn luyện $t$. \\
$\gamma$ & Hệ số chiết khấu trong cập nhật giá trị của học tăng cường. \\
$\tau$ & Hệ số cập nhật mềm của mạng đích. \\
$Q(s,a)$ & Giá trị hành động ước lượng bởi mạng Critic. \\
$V(s)$ & Giá trị trạng thái ước lượng bởi mạng giá trị của PPO. \\
\end{longtable}

\section{Số chiều trạng thái và hành động}
\label{sec:so-chieu}

Với cách ghép phần thực, phần ảo của ba khối kênh và chỉ báo phía người dùng, số chiều
trạng thái và số chiều hành động được xác định bởi:
\begin{equation}
\label{eq:a-so-chieu}
d_s = 2(K + N + KN) + K,
\qquad
d_a = (K+1) + K + 3N = 2K + 1 + 3N.
\end{equation}

Trong đó: $d_s$ là số phần tử của véc-tơ trạng thái; $d_a$ là số đầu ra chuẩn hóa của
chính sách; $K$ là số UE; $N$ là số phần tử STAR--RIS. Thành phần $2(K + N + KN)$ đến từ
phần thực và phần ảo của kênh trực tiếp, kênh trạm gốc--STAR--RIS và kênh STAR--RIS--UE;
thành phần $K$ cuối của $d_s$ cho biết mỗi UE nằm ở phía truyền hay phản xạ. Trong $d_a$,
$K+1$ phần tử dành cho công suất, $K$ phần tử dành cho phân bổ tốc độ chung và $3N$ phần
tử dành cho chia năng lượng cùng hai véc-tơ pha.

Công thức~\eqref{eq:a-so-chieu} cho thấy khi tăng số phần tử STAR--RIS, cả lượng thông tin
kênh mà mạng phải quan sát và số quyết định cần tạo đều tăng tuyến tính theo $N$. Đây là
lý do luận văn đánh giá nhiều giá trị $N$ để kiểm tra khả năng mở rộng của cả DRL và AO.

\begin{table}[htbp]
\centering
\caption{Số chiều trạng thái và hành động khi $K = 4$}
\label{tab:a-so-chieu}
\footnotesize
\begin{tabular}{@{}lrr@{}}
\toprule
$N$ & $d_s$ & $d_a$ \\
\midrule
$16$  & $172$  & $57$  \\
$32$  & $332$  & $105$ \\
$64$  & $652$  & $201$ \\
$96$  & $972$  & $297$ \\
$128$ & $1292$ & $393$ \\
\bottomrule
\end{tabular}
\tablesource{Tác giả tính từ cấu trúc trạng thái và hành động của mô hình (2026)}
\end{table}

Bảng~\ref{tab:a-so-chieu} giúp hình dung trực tiếp độ tăng của bài toán. Chẳng hạn, tại
$N = 32$, Actor nhận $332$ giá trị trạng thái và sinh $105$ giá trị hành động chuẩn hóa.
Các giá trị hành động này chưa phải biến vật lý cuối cùng mà tiếp tục đi qua bộ giải mã có
cấu trúc trước khi được dùng để tính kênh và tốc độ.

\chapter{QUY TRÌNH TÁI LẬP VÀ ĐỐI CHIẾU MÃ NGUỒN}
\label{app:tai-lap}

\section{Nguyên tắc tách dữ liệu}
\label{sec:tach-du-lieu}

Mỗi kích thước STAR--RIS sử dụng ba ngân hàng kịch bản kênh (scenario bank) độc lập cho huấn luyện, xác
thực và kiểm thử. Ngân hàng huấn luyện chứa $10.000$ kịch bản, ngân hàng xác thực chứa
$1.000$ kịch bản và ngân hàng kiểm thử chứa $1.000$ kịch bản. Thông tin mô tả và mã kiểm
tra SHA--256 được lưu cùng dữ liệu để phát hiện thay đổi ngoài ý muốn. Một bước kiểm tra
riêng xác nhận các ngân hàng không chứa cùng một kịch bản kênh.

Tập huấn luyện chỉ dùng để cập nhật tham số mạng. Tập xác thực được dùng định kỳ để lựa
chọn điểm lưu mô hình theo nguyên tắc ưu tiên QoS trước, tổng tốc độ sau. Tập kiểm thử
được dùng để báo cáo kết quả cuối sau khi quy tắc lựa chọn đã cố định. Khi so sánh phương
pháp tại cùng $N$, mọi phương pháp đều được đánh giá trên cùng các kịch bản kiểm thử để có
thể thực hiện so sánh ghép cặp theo trạng thái kênh.

\section{Thông tin cần lưu cho mỗi lần chạy}
\label{sec:thong-tin-luu}

Mỗi lần huấn luyện DRL cần lưu tối thiểu: tên thuật toán, hạt giống ngẫu nhiên, kích thước
STAR--RIS, số tương tác môi trường, tệp cấu hình, mã kiểm tra ngân hàng kênh, điểm lưu mô
hình được chọn, nhật ký huấn luyện, kết quả xác thực, kết quả kiểm thử, thời gian huấn
luyện và thông tin môi trường phần mềm cùng phần cứng. Việc lưu đồng thời các thành phần
này cho phép truy ngược một con số trong bảng kết quả về đúng lần chạy đã tạo ra nó.

Đối với các bộ giải xác định như AO--SCA, AO--Grid và AnalyticalRIS, không có biến thiên do
khởi tạo mạng. Tuy nhiên, vẫn cần lưu cấu hình bộ giải, số lần đánh giá hàm mục tiêu, tiêu
chí dừng, thời gian xử lý và danh sách kịch bản kiểm thử. Điều này đặc biệt quan trọng khi
so sánh độ trễ vì một thay đổi nhỏ về số vòng lặp hoặc kích thước lưới có thể làm thời gian
xử lý thay đổi đáng kể.

\section{Ánh xạ giữa nội dung luận văn và mã nguồn}
\label{sec:anh-xa-ma-nguon}

Bảng~\ref{tab:b-anh-xa} cho biết vị trí chính trong mã nguồn tương ứng với từng nội dung
của luận văn. Tên tệp được giữ theo dự án phần mềm để người đọc có thể đối chiếu trực tiếp.

Toàn bộ kết quả báo cáo trong Chương~\ref{chap:ket-qua} được sinh từ nhánh
\texttt{experiment/}\allowbreak\texttt{physical-v6-full}, mã phiên bản \texttt{fde2fe7},
với tham số hóa hành động \texttt{physical\_v6\_}\allowbreak\texttt{soft\_anchor}. Dữ liệu tổng hợp nằm tại thư mục
\texttt{results/physical\_v6\_full\_r2/}. Các nhánh khác của dự án chứa những lần chạy
trước đó và không dùng để trích dẫn.

\begin{table}[htbp]
\centering
\caption{Ánh xạ nội dung luận văn với các tệp mã nguồn chính}
\label{tab:b-anh-xa}
\footnotesize
\begin{tabularx}{\textwidth}{@{}X l@{}}
\toprule
\textbf{Nội dung} & \textbf{Tệp mã nguồn chính} \\
\midrule
Sinh kênh, hệ số STAR--RIS, kênh hiệu dụng, tốc độ RSMA & \texttt{src/star\_ris\_rsma/physics.py} \\
Trạng thái, chuẩn hóa, QoS và phần thưởng môi trường & \texttt{src/star\_ris\_rsma/env.py} \\
Bộ giải mã hành động có cấu trúc & \texttt{src/star\_ris\_rsma/action.py} \\
Kiến trúc Actor và Critic & \texttt{src/star\_ris\_rsma/agents/networks.py} \\
Cập nhật DDPG & \texttt{src/star\_ris\_rsma/agents/ddpg.py} \\
Cập nhật TD3 & \texttt{src/star\_ris\_rsma/agents/td3.py} \\
Cập nhật PPO & \texttt{src/star\_ris\_rsma/agents/ppo.py} \\
Bộ điều khiển hệ số phạt QoS thích nghi & \texttt{src/star\_ris\_rsma/experiment\_v2.py} \\
Huấn luyện chung, xác thực và chọn điểm lưu mô hình & \texttt{src/star\_ris\_rsma/experiment\_v3.py} \\
Ngân hàng kịch bản kênh và mã kiểm tra & \texttt{src/star\_ris\_rsma/scenario\_bank.py} \\
AO--SCA & \texttt{src/star\_ris\_rsma/baselines/ao\_corrected.py} \\
AO--Grid & \texttt{src/star\_ris\_rsma/baselines/ao\_grid\_corrected.py} \\
Các phép loại bỏ thành phần để phân tích & \texttt{src/star\_ris\_rsma/baselines/ablations.py} \\
\bottomrule
\end{tabularx}
\tablesource{Tác giả tổng hợp từ cấu trúc dự án phần mềm (2026)}
\end{table}

Các hậu tố trong một số tên tệp là ký hiệu kỹ thuật nội bộ của dự án phần mềm. Ví dụ gồm
\texttt{experiment\_v2.py}, \texttt{experiment\_v3.py} và \texttt{ao\_corrected.py}. Các
ký hiệu này không biểu thị các phương pháp khoa học khác nhau trong nội dung luận văn.
Luận văn chỉ mô tả phương pháp cuối cùng tương ứng với cấu hình và kết quả đã khóa.

\section{Danh sách kiểm tra trước khi sử dụng kết quả}
\label{sec:danh-sach-kiem-tra}

Trước khi đưa một bảng hoặc đồ thị vào luận văn, cần thực hiện các bước sau:

\begin{enumerate}[leftmargin=*,itemsep=2pt]
\item Xác nhận ba ngân hàng huấn luyện, xác thực và kiểm thử có đúng số lượng, đúng kích
      thước $K$ và $N$, đồng thời không trùng kịch bản.
\item Xác nhận mã kiểm tra của ngân hàng kênh khớp tệp kê khai đã khóa.
\item Xác nhận đủ năm hạt giống huấn luyện cho mỗi thuật toán DRL và mỗi $N$; không loại
      một lần chạy chỉ vì kết quả thấp.
\item Kiểm tra các tệp kết quả thô không chứa giá trị không xác định hoặc vô hạn.
\item Xác nhận DDPG, TD3, PPO, AO--SCA, AO--Grid và AnalyticalRIS được đánh giá trên cùng
      $1.000$ kịch bản kiểm thử tại mỗi $N$.
\item Xác nhận các bảng tổng hợp có thể tái tạo từ dữ liệu thô, không nhập tay số liệu
      trung gian.
\item Khi báo cáo độ trễ, xác nhận các phương pháp được đo trên cùng máy, cùng số luồng xử
      lý, cùng số lần khởi động và số lần lặp đo.
\item Khi báo cáo kiểm định thống kê, phân biệt rõ đơn vị kịch bản với đơn vị hạt giống
      huấn luyện; không dùng kết quả theo kịch bản để thay thế biến thiên do huấn luyện.
\item Kiểm tra mọi công thức và giá trị siêu tham số trong luận văn khớp mã nguồn và cấu
      hình cuối cùng.
\item Kiểm tra mọi hình và bảng có nguồn, chú thích, đơn vị và mô tả đủ để đọc độc lập với
      phần mã nguồn.
\end{enumerate}

\section{Nguyên tắc đọc kết quả thống kê}
\label{sec:nguyen-tac-thong-ke}

Đối với DRL, một kịch bản kênh có thể được đánh giá bởi nhiều chính sách đến từ các hạt
giống huấn luyện khác nhau. Trung bình theo kịch bản cho biết các chính sách đã cố định
hoạt động thế nào trên các trạng thái kênh chung, còn trung bình theo hạt giống phản ánh
sự biến thiên của quá trình huấn luyện. Hai đơn vị không thể thay thế cho nhau. Khi số hạt
giống chỉ bằng năm, luận văn ưu tiên các nhận định về xu hướng và độ lớn hiệu ứng, đồng
thời tránh kết luận mạnh chỉ dựa trên giá trị xác suất nhỏ từ $1.000$ kịch bản ghép cặp.

Đối với AO--SCA, AO--Grid và AnalyticalRIS, thuật toán là xác định nên không có độ lệch
chuẩn theo hạt giống huấn luyện. Khoảng biến thiên của chúng chủ yếu phản ánh độ khó khác
nhau giữa các kịch bản kênh. Vì vậy, độ rộng khoảng tin cậy của AO không được so trực tiếp
với độ lệch chuẩn theo hạt giống của DRL như thể hai đại lượng cùng biểu diễn một nguồn
bất định.
